% Môn: Vật lý
% Lớp: 10
% Chương: Chương VII. Biến dạng của vật rắn. Áp suất chất lỏng
% Bài: L10C7 Bài 33. Biến dạng của vật rắn
% Số câu: 7
% App đọc file này trực tiếp — sửa rồi Commit trên GitHub, bấm Đồng bộ GitHub .tex

% L10C7 Bài 33. Biến dạng của vật rắn

% ===== Câu 1 =====
% ID: L10C7B33-01-DS
% Mức: NB
\begin{ex}
% ID: L10C7B33-01-DS
% Mức: NB
Một chuyển động thẳng biến đổi đều có
\choiceTF
{\True véc tơ gia tốc không đổi}
{\True gia tốc của chuyển động không đổi}
{vận tốc của chuyển động tăng dần đều theo thời gian}
{\True vận tốc của chuyển động là hàm bậc nhất của thời gian}
\loigiai{
\begin{itemchoice}
\item (Đúng) véc tơ gia tốc không đổi. (Vì): Trong chuyển động thẳng, phương của $\vec a$ cố định; “biến đổi đều” $\Rightarrow\vec a$ không đổi
\item (Đúng) gia tốc của chuyển động không đổi. (Vì): ây là phát biểu trực tiếp định nghĩa
\item (Sai) vận tốc của chuyển động tăng dần đều theo thời gian. (Vì): Nếu là chậm dần đều thì tốc độ giảm đều, không tăng
\item (Đúng) vận tốc của chuyển động là hàm bậc nhất của thời gian. (Vì): $v=v_0+at$ là hàm bậc nhất của $t$ khi $a$ không đổi
\end{itemchoice}
}
\end{ex}

% ===== Câu 2 =====
% ID: L10C7B33-01-TN
% Mức: NB
\begin{ex}
% ID: L10C7B33-01-TN
% Mức: NB
Câu nào sau đây sai về lực căng?
\choice
{Lực căng có phương trùng với chính sợi dây, chiều hướng từ hai đầu vào phần giữa của dây}
{Lực căng của dây có bản chất là lực đàn hồi}
{\True Lực căng có thể là lực kéo hoặc lực nén}
{Lực căng của dây có điểm đặt là điểm mà đầu dây tiếp xúc với vật}
\loigiai{
(1) Đề bài yêu cầu xác định phát biểu sai về lực căng.
(2) Phân tích từng phương án:
- A. Lực căng có phương trùng với chính sợi dây, chiều hướng từ hai đầu vào phần giữa của dây. Đúng, đây là đặc điểm của lực căng trong dây bị kéo căng.
- B. Lực căng của dây có bản chất là lực đàn hồi. Đúng, lực căng là một dạng của lực đàn hồi xuất hiện khi dây bị kéo căng.
- C. Lực căng có thể là lực kéo hoặc lực nén. Sai, lực căng chỉ có thể là lực kéo, không thể là lực nén. Dây chỉ có thể bị kéo căng, không thể chịu lực nén mà vẫn giữ được tính chất căng.
- D. Lực căng của dây có điểm đặt là điểm mà đầu dây tiếp xúc với vật. Đúng, lực căng tác dụng lên vật tại điểm tiếp xúc.
(3) Phát biểu sai là C.
}
\end{ex}

% ===== Câu 3 =====
% ID: L10C7B33-02-DS
% Mức: NB
\begin{ex}
% ID: L10C7B33-02-DS
% Mức: NB
Vectơ gia tốc của chuyển động thẳng biến đổi đều
\choiceTF
{có phương vuông góc với vectơ vận tốc}
{\True có độ lớn không đổi}
{cùng hướng với vectơ vận tốc}
{ngược hướng với vectơ vận tốc}
\loigiai{
\begin{itemchoice}
\item (Sai) có phương vuông góc với vectơ vận tốc. (Vì): Chuyển động thẳng nên $\vec a$ luôn cùng phương với $\vec v$, không vuông góc
\item (Đúng) có độ lớn không đổi. (Vì): “Biến đổi đều” nghĩa là gia tốc có độ lớn không đổi
\item (Sai) cùng hướng với vectơ vận tốc. (Vì): Chỉ đúng với chuyển động thẳng \emph{nhanh dần đều}. Với chậm dần đều thì ngược hướng
\item (Sai) ngược hướng với vectơ vận tốc. (Vì): Chỉ đúng với chuyển động thẳng \emph{chậm dần đều}. Với nhanh dần đều thì cùng hướng
\end{itemchoice}
}
\end{ex}

% ===== Câu 4 =====
% ID: L10C7B33-02-TN
% Mức: NB
\begin{ex}
% ID: L10C7B33-02-TN
% Mức: NB
Kết luận nào sau đây không đúng? Đối với lực đàn hồi của lò xo
\choice
{xuất hiện khi vật bị biến dạng}
{ngược hướng với lực làm nó bị biến dạng}
{tỉ lệ với độ biến dạng}
{\True luôn là lực kéo}
\loigiai{
**Phân tích các phương án:**
 * **A. xuất hiện khi vật bị biến dạng**: Đúng. Lực đàn hồi là lực phục hồi, chỉ xuất hiện khi vật bị biến dạng (kéo giãn hoặc nén).
 * **B. ngược hướng với lực làm nó bị biến dạng**: Đúng. Lực đàn hồi luôn có xu hướng chống lại sự biến dạng, đưa vật về trạng thái ban đầu. Ví dụ, khi kéo lò xo xuống, lực đàn hồi hướng lên; khi nén lò xo vào, lực đàn hồi đẩy ra.
 * **C. tỉ lệ với độ biến dạng**: Đúng. Theo định luật Hooke, độ lớn lực đàn hồi tỉ lệ thuận với độ lớn độ biến dạng của lò xo ($F_{đh} = k \cdot |\Delta l|$), trong giới hạn đàn hồi.
 * **D. luôn là lực kéo**: Sai. Lực đàn hồi có thể là lực kéo (k
}
\end{ex}

% ===== Câu 5 =====
% ID: L10C7B33-03-DS
% Mức: NB
\begin{ex}
% ID: L10C7B33-03-DS
% Mức: NB
Chọn ý Khi một chất điểm chuyển động thẳng biến đổi đều thì nó có
\choiceTF
{\True gia tốc không đổi}
{\True tốc độ tức thời tăng đều hoặc giảm đều theo thời gian}
{gia tốc tăng dần đều theo thời gian}
{\True có thể lúc đầu chậm dần đều, sau đó nhanh dần đều}
\loigiai{
\begin{itemchoice}
\item (Đúng) gia tốc không đổi. (Vì): ịnh nghĩa “biến đổi đều” là $a$ không đổi
\item (Đúng) tốc độ tức thời tăng đều hoặc giảm đều theo thời gian. (Vì): Khi $a$ không đổi, tốc độ tăng đều (nhanh dần) hoặc giảm đều
\item (Sai) gia tốc tăng dần đều theo thời gian. (Vì): Gia tốc không đổi nên không “tăng dần” theo thời gian
\item (Đúng) có thể lúc đầu chậm dần đều, sau đó nhanh dần đều.
\end{itemchoice}
}
\end{ex}

% ===== Câu 6 =====
% ID: L10C7B33-04-DS
% Mức: NB
\begin{ex}
% ID: L10C7B33-04-DS
% Mức: NB
Cho ba xe (1), (2) và (3) chuyển động thẳng (đồ thị $x(t)$ như Hình. Hãy xác định đúng–sai cho các phát biểu sau:
\choiceTF
{\True Xe (1) đứng yên so với xe (3).}
{Vận tốc tương đối của xe (1) đối với xe (2) bằng vận tốc tương đối của xe (2) đối với xe (3).}
{\True Người ngồi trên xe (2) thấy xe (1) và xe (3) cùng chạy về phía mình (xe (1) ở phía trước, xe (3) phía sau).}
{\True Người ngồi trên xe (3) thấy xe (1) và xe (2) chạy ra xa mình.}
\loigiai{
\begin{itemchoice}
\item (Đúng) Xe (1) đứng yên so với xe (3). (Vì): Từ đồ thị: (1) và (3) có cùng độ dốc (vận tốc như nhau) nên đứng yên tương đối $\Rightarrow$
\item (Sai) Vận tốc tương đối của xe (1) đối với xe (2) bằng vận tốc tương đối của xe (2) đối với xe (3). (Vì): Vận tốc tương đối không có tính “bắc cầu” $v_{1/2}\neq v_{2/3}$ nói chung $\Rightarrow$
\item (Đúng) Người ngồi trên xe (2) thấy xe (1) và xe (3) cùng chạy về phía mình (xe (1) ở phía trước, xe (3) phía sau).
\item (Đúng) Người ngồi trên xe (3) thấy xe (1) và xe (2) chạy ra xa mình.
\end{itemchoice}
}
\end{ex}

% ===== Câu 7 =====
% ID: L10C7B33-05-DS
% Mức: NB
\begin{ex}
% ID: L10C7B33-05-DS
% Mức: NB
Gia tốc là một đại lượng
\choiceTF
{đại số, đặc trưng cho sự biến thiên nhanh hay chậm của chuyển động}
{đại số, đặc trưng cho tính không đổi của vận tốc}
{vectơ, đặc trưng cho sự biến thiên nhanh hay chậm của chuyển động}
{\True vectơ, đặc trưng cho sự biến thiên nhanh hay chậm của vận tốc}
\loigiai{
\begin{itemchoice}
\item (Sai) đại số, đặc trưng cho sự biến thiên nhanh hay chậm của chuyển động. (Vì): Gia tốc gắn với sự \emph{biến thiên vận tốc}, không chỉ chung chung “chuyển động” và không phải đại lượng vô hướng
\item (Sai) đại số, đặc trưng cho tính không đổi của vận tốc. (Vì): “Tính không đổi của vận tốc” thì gia tốc phải bằng $0$
\item (Sai) vectơ, đặc trưng cho sự biến thiên nhanh hay chậm của chuyển động. (Vì): Diễn đạt “chuyển động nhanh, chậm” chưa chuẩn; về bản chất gia tốc đặc trưng cho sự thay đổi của \emph{vận tốc}
\item (Đúng) vectơ, đặc trưng cho sự biến thiên nhanh hay chậm của vận tốc. (Vì): Gia tốc là một đại lượng vectơ, đặc trưng cho sự biến thiên nhanh hay chậm của vận tốc theo thời gian
\end{itemchoice}
}
\end{ex}
